% !TEX root = ../main.tex
\clearpage
\phantomsection
\section*{4 \textbf{이론적 틀}}
\label{sec:theoretical-framework}
\addcontentsline{toc}{section}{4 이론적 틀}

본 연구는 하이퍼링크-인용 이론, 구성개념(도메인) 정렬, 계산 복잡도 이론에 기반하여 EndorseRank를 Adaptive Weighted PageRank(AWP)와 세 가지 핵심 구성개념---평판 구조, 거래 성공 정렬, 계산 효율성---에 걸쳐 비교한다. 각 구성개념은 앞서 연구 목표 및 기대 결과 틀에서 소개한 파라미터 정의와 직접 연결되며 검증 가능한 가설을 지원한다.


\phantomsection
\subsection*{4.1 구성개념 및 변수 매핑}
\label{sec:constructs-and-variable-mapping}
\addcontentsline{toc}{subsection}{4.1 구성개념 및 변수 매핑}

표~1은 각 평가 구성개념을 변수 및 운영 측정치에 매핑한다.

\begin{longtable}{>{\raggedright\arraybackslash}p{0.25\linewidth}>{\raggedright\arraybackslash}p{0.36\linewidth}>{\raggedright\arraybackslash}p{0.29\linewidth}}
\caption{평가 구성개념 및 관련 변수}\\
\toprule
\textbf{구성개념} & \textbf{관련 변수} & \textbf{운영 측정치} \\
\midrule
평판 구조 & EndorseRank 점수, AWP 점수, 순위 순서, 순위 발산 & Spearman's rho, Kendall's tau, top-k 겹침 \\
\addlinespace
거래 성공 정렬 & GMX 청산 성공 횟수, 실현 이익 대리변수, 청산 성공률, 평판 점수 & 각 대리변수 대비 Spearman's rho, Kendall's tau \\
\addlinespace
계산 효율성 & 그래프 구축 시간, PageRank 수렴 시간, 메모리 사용량, 반복 횟수 & 런타임, 최대 메모리, 허용오차까지 반복 \\
\addlinespace
\bottomrule
\end{longtable}

\phantomsection
\subsection*{4.2 개념 모형}
\label{sec:conceptual-model}
\addcontentsline{toc}{subsection}{4.2 개념 모형}

그림~1은 개념 모형을 제시한다: EndorseRank 대 AWP 비교는 세 구성개념(평판 구조, 거래 성공 정렬, 계산 효율성)에 걸쳐 평가되며, 각각 해당 검증 가능 가설(H1--H3)을 갖는다.


\begin{center}
% !TEX root = ../main.tex
\begin{tikzpicture}[
  font=\small,
  box/.style={draw, rounded corners, align=center, minimum width=10.8cm, minimum height=1.0cm},
  smallbox/.style={draw, rounded corners, align=center, minimum width=3.4cm, minimum height=1.05cm},
  arrow/.style={->, thick}
]
\node[box] (center) at (0,2.6) {EndorseRank vs. AWP};

\node[smallbox] (rep) at (-5,0.6) {평판\\구조};
\node[smallbox] (risk) at (0,0.6) {거래 성공\\정렬};
\node[smallbox] (eff) at (5,0.6) {계산\\효율성};

\node[smallbox] (h1) at (-5,-1.4) {H1: Divergence};
\node[smallbox] (h2) at (0,-1.4) {H2: 거래 성공\\정렬};
\node[smallbox] (h3) at (5,-1.4) {H3: 계산\\효율성};

\draw[arrow] (center.south) -- (rep.north);
\draw[arrow] (center.south) -- (risk.north);
\draw[arrow] (center.south) -- (eff.north);

\draw[arrow] (rep.south) -- (h1.north);
\draw[arrow] (risk.south) -- (h2.north);
\draw[arrow] (eff.south) -- (h3.north);
\end{tikzpicture}


\vspace{0.4em}
그림 1: EndorseRank vs. AWP: 가설 및 평가 구성개념
\end{center}

\textbf{H1 (Divergence): }EndorseRank와 AWP는 상당히 다른 지갑 순서를 산출하여, approve/permit 기반 보증 그래프가 구별되는 평판 구조를 만든다.


\textbf{H2 (거래 성공 정렬): }EndorseRank 점수는 GMX V2 거래 성공 대리변수(청산 성공 횟수 및 실현 이익 대리변수)와 AWP 점수보다 더 강하게 상관하여, 우수한 도메인 직관 정렬을 나타낸다.


\textbf{H3 (계산 효율성): }EndorseRank는 AWP보다 더 적은 런타임, 메모리, 수렴 반복을 요구하여 실용적 확장성 이점을 입증한다.


\phantomsection
\subsection*{4.3 이론적 정당화}
\label{sec:theoretical-justification}
\addcontentsline{toc}{subsection}{4.3 이론적 정당화}

\textbf{하이퍼링크--인용 의미론: }원래 PageRank 알고리즘(Page et al., 1998)은 들어오는 하이퍼링크 구조로 웹페이지의 중요성을 평가하며, 각 링크를 인용으로 취급하고 감쇠 확률로 링크를 따르는 무작위 서퍼를 모델링한다. Do et al.(2023)은 이 전파 논리를 이더리움 거래 그래프에 적용하여 엣지를 송금 가치와 시간 감쇠로 가중하여 평판이 관찰된 경제적 연결을 따라 흐르도록 한다. 본 연구는 승인(allowance)이 명시적 온체인 지출 권한 위임이기 때문에(Ethereum Improvement Proposals, 2015; Ethereum Improvement Proposals, 2020) 동일한 인용 기반 전파를 ERC-20 approve/permit 엣지로 확장한다. 토큰 송금과 달리 거래나 자동 라우팅을 반영할 수 있으나 신뢰 의도 없이 발생할 수 있는 반면, 승인은 취소될 때까지 지속적인 승인 상태를 기록한다. 이더리움의 그래프 기반 리스크 모델(Lin et al., 2024)은 네트워크 구조가 원시 활동 수를 넘어 리스크 관련 정보를 담고 있음을 추가로 보여주며, 송금량만이 아닌 방향성 보증 그래프 사용을 지지한다.


\textbf{도메인 정렬 프레임워크: }측정 이론에서 구성개념 정렬은 지표가 관심 결과를 얼마나 잘 추적하는지 평가한다(Cronbach \& Meehl, 1955). 순위 기반 평판 점수 산출에 Spearman's rho와 Kendall's tau가 적합한 이유는 분류 임계값을 가정하지 않고 평판 점수와 도메인 대리변수 간 단조적·쌍별 순서 일치를 평가하기 때문이다(Do et al., 2023). Do et al.(2023)이 PageRank를 in-degree 및 in-value에 대해 검증한 것을 따르면, GMX V2 \textbf{청산 성공 횟수}와 \textbf{실현 이익 대리변수}를 Arbitrum One에서 관찰 가능한 거래 성공 측정치로 취급한다. 이러한 대리변수와 정렬하는 평판 점수는 외부 ground-truth 대출 상환 레이블을 주장하지 않고 본 연구의 도메인 직관 요건을 충족한다.


\textbf{레버리지 노출 및 거래 성공 정렬: }GMX V2 무기한 스왑 포지션은 담보화된 레버리지 노출을 나타낸다: 트레이더는 마진을 예치하고 롱 또는 숏 포지션을 통해 기초 자산에 대한 방향성 리스크를 부담한다. 수익성 \textbf{PositionDecrease} 청산(양의 \textbf{basePnlUsd}, 포지션 청산 제외)은 차입 또는 레버리지 노출을 성공적으로 관리하는 것---신용 맥락에서 차입과 의무 이행에 유사---과 기능적으로 닮았다.


\textbf{계산 복잡도 이론: }PageRank의 반복당 시간 복잡도는 $O(|E|)$이다(Page et al., 1998). approve/permit 그래프는 송금 기반 그래프보다 희소(엣지 $|E|$가 더 적음)하여 반복당 비용 감소, 더 빠른 수렴, 메모리 사용량 감소로 이어진다.


\textbf{범위 및 향후 확장: }본 계획서는 블랙박스 예측 모델보다 해석 가능한 그래프 전파와 검증 가능한 온체인 신호를 우선한다. DeFi 신용 점수에 대한 지도 학습 및 표현 학습 접근(Ghosh et al., 2024; Kotb, 2023)은 순위 성능을 개선할 수 있으나 Packin and Lev-Aretz(2024)가 강조한 가명성 제약 하에서 투명성과 감사 가능성을 희생하는 경우가 많다. EndorseRank는 따라서 각 점수가 명시적 승인 엣지에 추적 가능하도록 PageRank 계열 내에 머문다. 향후 연구는 보증 그래프 특성을 DeFi 정보 표현 학습과 결합할 수 있으나, 그러한 확장은 본 연구 범위 밖이다.


\phantomsection
\subsection*{4.4 연구 설계에 대한 함의}
\label{sec:implications-for-research-design}
\addcontentsline{toc}{subsection}{4.4 연구 설계에 대한 함의}

연구 설계는 다음 측정 및 분석 단계를 통해 위 구성개념을 적용한다.

\textbf{측정: }Arbitrum One에서 GMX V2 활성 지갑의 매칭 표본에 대해 EndorseRank 및 AWP 점수를 계산하고, 정의된 관찰 기간 동안 \textbf{PositionDecrease} 이벤트에서 청산 성공 횟수, 실현 이익 대리변수, 청산 성공률을 도출하며, 각 알고리즘의 런타임, 메모리 사용량, 수렴까지 반복 횟수를 프로파일링한다.


가설 검증 및 해석은 위 측정 단계를 따르며 다음과 같이 구성된다.


\proposalSubheading{분석:}

\textbf{H1 }은 EndorseRank와 AWP 점수 간 순위 일치(Spearman's \emph{\(\rho\)}, Kendall's \emph{\(\tau\)}, top-\emph{k} 겹침)를 통해, 이러한 계수가 완전 일치에서 유의미하게 다른지 검증


\textbf{H2 }는 평판 점수와 GMX 거래 성공 대리변수 순위 간 순위 상관(Spearman's \emph{\(\rho\)}, Kendall's \emph{\(\tau\)})을 통해. Spearman's rho는 단조적 순위 일치를, Kendall's tau는 쌍별 순서 일치를 측정한다. 이 지표는 평판 점수와 거래 성공 대리변수가 이진 부도 레이블이 아닌 연속 순위 변수이므로 정당화되며, Do et al.(2023)의 도메인 정렬 접근을 따른다.


\textbf{H3 }은 런타임, 최대 메모리, 반복 횟수의 쌍별 벤치마크를 통해


\textbf{기여: }가설을 확립된 이론에 고정함으로써, 이 틀은 실증 결과가 학술적 엄밀성과 실무적 관련성을 모두 제공하여 EndorseRank가 의미론적으로 유의미하고 도메인에 정렬되며 효율적인 온체인 평판 지표를 산출하는 이점을 입증할 것임을 보장한다.
